% Owner: Abhinav (P8.5); details in docs/data_sources.md
\section{Data}\label{sec:data}

All data are free and openly licensed (Table~\ref{tab:data}). Rasters are read from the Microsoft
Planetary Computer, and only the window needed for the study area plus a \SI{3}{\km} buffer is
downloaded. The buffer equals the distance cap of the features, so distances from cells near the
edge are exact. OpenStreetMap layers come from the Overpass API through OSMnx \citep{boeing2017}.
Every download is recorded with its checksum and extent in a manifest, so a rerun reproduces or
skips it.

\begin{table}[ht]
  \centering\small
  \caption{Datasets.}
  \label{tab:data}
  \begin{tabularx}{\textwidth}{@{}X l l l X@{}}
    \toprule
    Dataset & Years & Resolution & Licence & Used for \\
    \midrule
    ESRI 10\,m Annual LULC v02 \citep{karra2021} & 2018--2023 & \SI{10}{\m} & CC BY 4.0 &
      land-cover features, three-class map, growth labels \\
    ESA WorldCover v200 \citep{zanaga2022} & 2021 & \SI{10}{\m} & CC BY 4.0 &
      cross-check of the three-class map \\
    Copernicus DEM GLO-30 \citep{copdem} & static & \SI{30}{\m} & Copernicus DEM licence &
      elevation, slope \\
    OpenStreetMap \citep{osm} & 2018-01-01, current & vector & ODbL 1.0 &
      roads, water, buildings, protected areas \\
    \bottomrule
  \end{tabularx}
\end{table}

\paragraph{Land use / land cover.} The ESRI (Impact Observatory) map is used for every year and
every purpose, so that ``built'' has the same definition in the baseline and the test years. It
has the highest overall accuracy of the global 10\,m products (75\,\%), and built area is one of
its most reliable classes \citep{venter2022}. Its known bias, over-estimating scrub, matters
here: much of the Bannerghatta forest appears as rangeland. Inside the park, ESRI classes 47\,\%
of the land as trees and 50\,\% as rangeland. ESA WorldCover 2021 is used only as an independent
cross-check. The two products disagree strongly on built-up land (62\,\% of the area in ESRI 2023
vs 33\,\% in WorldCover), because ESRI generalises urban areas to include gardens and street
trees. ESRI land cover also flickers between years: only 35--49\,\% of single-year ``new built''
pixels are still built later. The growth labels therefore require two years at each end
(Section~\ref{sec:validation_design}).

\paragraph{Terrain.} The Copernicus DEM is warped to UTM at \SI{30}{\m}. Slope is computed with
Horn's method \citep{horn1981} and then resampled to the \SI{10}{\m} grid. The DEM is a surface
model, so heights include buildings and trees; slope is less reliable in dense built-up cells, but
it matters most on open candidate land, where this effect is small.

\paragraph{OpenStreetMap.} Roads, water and buildings are downloaded twice: as they were on
1 January 2018 (an Overpass \texttt{[date:]} query) for the validation run, and as they are now
for the final map. The difference is large. All-road length more than doubled between the two
snapshots (\SI{3889}{\km} to \SI{8730}{\km}), partly new streets in growth areas and partly later
mapping. Road features from today's map would therefore leak future growth into a 2018 analysis.
OSM building coverage is also very uneven: buildings were mapped in 62--75\,\% of the built cells
in the north but in about 1\,\% in the south in 2018, and almost all are untagged
(\texttt{building=yes}). Protected areas (Bannerghatta, tagged as a national park) are a static
layer and come from the current map.

\paragraph{Preprocessing.} All rasters are aligned to one reference grid in EPSG:32643: \SI{10}{\m}
pixels whose edges fall on whole \SI{100}{\m} cells, so every block of $10\times10$ pixels is one
grid cell. Land cover is copied without resampling, because the ESRI tile is already on this
lattice. OSM layers are reprojected, clipped to the grid, repaired and split into single-part
geometries. Drivable roads are split into major (motorway to secondary) and minor. Buildings
smaller than \SI{10}{\m\squared} or larger than \SI{10}{\hectare} are dropped. Each preparation
step is checked independently against the source data, for example by comparing pixel values at
random points, matching road lengths before and after clipping, and overlaying OSM lakes on ESRI
water.
